% Patrones de flujo en una tolva y tolva en cuña sobre el tornillo.
% Elaboración propia con base en Su (2017), Mehos y Morgan (2016) y Bates (2000).
% Variables de las Ecuaciones de fricción de pared y abertura mínima: theta y B.
\begin{tikzpicture}[
    >=Latex,
    font=\small,
    pared/.style={thick},
    mov/.style={-{Latex[width=4pt,length=3pt]}, thin},
    cota/.style={<->, thin},
    nota/.style={font=\footnotesize, align=center}
]

% (a) Flujo másico: todo el contenido se mueve y desliza en las paredes
\begin{scope}
    \fill[gray!10] (0,4.4) -- (0,2.2) -- (1.2,0.5) -- (1.8,0.5) -- (3,2.2) -- (3,4.4) -- cycle;
    \draw[pared] (0,4.6) -- (0,2.2) -- (1.2,0.5) (1.8,0.5) -- (3,2.2) -- (3,4.6);
    \foreach \x in {0.35,1.5,2.65}{\draw[mov] (\x,4.1) -- (\x,3.2);}
    \draw[mov] (0.45,2.0) -- (1.05,1.15);
    \draw[mov] (2.55,2.0) -- (1.95,1.15);
    \draw[mov] (1.5,2.4) -- (1.5,1.2);
    \node[nota] at (1.5,-0.2) {(a) Flujo másico};
\end{scope}

% (b) Flujo en embudo: canal central en movimiento y zonas en reposo
\begin{scope}[xshift=4.2cm]
    \fill[gray!10] (0,4.4) -- (0,2.2) -- (1.2,0.5) -- (1.8,0.5) -- (3,2.2) -- (3,4.4) -- cycle;
    \fill[gray!45] (0,3.6) -- (0,2.2) -- (1.2,0.5) -- (1.15,1.1) -- cycle;
    \fill[gray!45] (3,3.6) -- (3,2.2) -- (1.8,0.5) -- (1.85,1.1) -- cycle;
    \draw[pared] (0,4.6) -- (0,2.2) -- (1.2,0.5) (1.8,0.5) -- (3,2.2) -- (3,4.6);
    \draw[thin, dashed] (1.15,1.1) -- (0,3.6) (1.85,1.1) -- (3,3.6);
    \foreach \x in {1.2,1.5,1.8}{\draw[mov] (\x,3.4) -- (\x,1.4);}
    \draw[thin] (0.3,2.3) -- (-0.25,2.95);
    \node[nota, anchor=east] at (-0.25,3.0) {Reposo};
    \node[nota] at (1.5,-0.2) {(b) Flujo en embudo};
\end{scope}

% (c) Tolva en cuña con salida de ranura sobre el tornillo
\begin{scope}[xshift=8.4cm]
    \fill[gray!10] (0,3.8) -- (1.1,0.9) -- (1.9,0.9) -- (3,3.8) -- cycle;
    \draw[pared] (0,3.8) -- (1.1,0.9) (1.9,0.9) -- (3,3.8);
    \draw[pared] (1.1,0.9) -- (1.1,0.5) arc[start angle=180, end angle=360, radius=0.4] -- (1.9,0.9);
    \draw[thick, fill=gray!30] (1.5,0.5) circle[radius=0.33];
    \draw[fill=white] (1.5,0.5) circle[radius=0.09];
    \node[nota, anchor=west] at (2.05,0.3) {Tornillo};
    % semiángulo theta, medido desde la vertical
    \draw[thin, dashed] (1.1,0.9) -- (1.1,2.6);
    \draw[thin] (1.1,1.9) arc[start angle=90, end angle=110.8, radius=1.0];
    \node[nota] at (0.88,2.2) {$\theta$};
    % ancho de la salida B
    \draw[thin, gray] (1.1,0.9) -- (1.1,1.0) (1.9,0.9) -- (1.9,1.0);
    \draw[cota] (1.1,1.12) -- (1.9,1.12);
    \node[nota] at (1.6,1.38) {$B$};
    \node[nota] at (1.5,-0.2) {(c) Tolva en cuña};
\end{scope}

\end{tikzpicture}
